\section*{Chapter \ref{ch:nw:long-haul-wireless} --- Long-Haul Wireless}

\begin{solution}{exo:nw:long-haul-wireless:1}
Bulge $20 \times 20 / (12.74 \times 4/3) = \qty{23.5}{\metre}$; Fresnel radius $17.32\sqrt{400/(11 \times 40)} = \qty{16.5}{\metre}$.
\end{solution}

\begin{solution}{exo:nw:long-haul-wireless:2}
$0.0007056 \times 25^{1.59} = \qty{0.118}{\decibel\per\kilo\metre}$ at \qty{6}{\giga\hertz}; $0.07078 \times 25^{1.0818} = \qty{2.30}{\decibel\per\kilo\metre}$ at
\qty{18}{\giga\hertz}, twenty times more.
\end{solution}

\begin{solution}{exo:nw:long-haul-wireless:3}
Towers can be placed almost anywhere a lease can be had, in a nearly straight line, and the signal crosses fields, rivers and roads without
needing a right of way along them; a cable must follow continuous corridors in the ground. The \omterm{def:nw:the-north-american-map:factor}{route factor} measures only the path; the air's
speed is a separate gain of 46\%.
\end{solution}

\begin{solution}{exo:nw:long-haul-wireless:4}
$20\log_{10} 2 = \qty{6.02}{\decibel}$ more loss, so \qty{6}{\decibel} less margin; and the rain acts over twice the length. In the model a 59-km
hop at \qty{11}{\giga\hertz} breaks at \qty{18.7}{\milli\metre\per\hour} of rain along its length; a 118-km hop at \qty{9.1}{\milli\metre\per\hour}.
\end{solution}

\begin{solution}{exo:nw:long-haul-wireless:5}
Both carry gigabits over the short hops of a city; they fail in different weather (millimetre waves in heavy rain, lasers in fog), so carrying
the traffic on both keeps the route up unless both kinds of weather come together.
\end{solution}

\begin{solution}{exo:nw:long-haul-wireless:6}
It can carry a few compact messages (an index future's best price, a signal) faster than any cable; it cannot carry a market-data feed, and it
is available only when the ionosphere allows, so it needs a cable behind it.
\end{solution}

\begin{solution}{exo:nw:long-haul-wireless:7}
Scanning the peak rate, the first storm that puts the route on fibre has a peak of \qty{48}{\milli\metre\per\hour}.
\end{solution}

\begin{solution}{exo:nw:long-haul-wireless:8}
Availabilities multiply along a chain: $0.9995^{20} = 0.990$, if the hops fail independently. The route is down about 1\% of the year, some
87 hours, twenty times more than one hop; and hops are not independent, since one storm can cover two neighbours, which makes the arithmetic
more complicated but does not restore 99.95\%.
\end{solution}

\begin{solution}{pb:nw:long-haul-wireless:1}
\textbf{Part I.}
\begin{enumerate}
\item \qty{51.3}{\metre} of bulge and \qty{20.1}{\metre} of Fresnel radius at mid-hop; towers of \qty{81.4}{\metre}.
\item \qty{148.7}{\decibel} of free-space loss, $-33.5$~dBm received, \qty{36.5}{\decibel} of \omterm{def:nw:long-haul-wireless:budget}{fade margin}.
\item \qty{18.7}{\milli\metre\per\hour}.
\item \qty{80.6}{\milli\metre\per\hour}.
\end{enumerate}
\textbf{Part II.}
\begin{enumerate}[resume]
\item $30/0.8 = 37.5$ minutes: the model counts 38 one-minute steps.
\item 21 minutes.
\item $6.911 - 3.982 = \qty{2.93}{\milli\second}$ one way.
\item At \qty{6}{\giga\hertz} the route never leaves the radio; at \qty{18}{\giga\hertz} it spends 29 minutes on fibre.
\end{enumerate}
\textbf{Part III.}
\begin{enumerate}[resume]
\item $0.05\,e^{-18.7/4} = 0.047\%$ of the year: about 4.1 hours.
\item $1 - (1 - 0.00047)^{20} = 0.94\%$: about 82 hours a year on fibre.
\item Rain is spatially correlated: one storm often covers several neighbouring hops, and a climate's wet season hits them all; failures cluster in
  time and space.
\item Less free-space loss and a shorter wet length per hop, so each hop tolerates heavier rain; but more towers, more radios and more
  microseconds of equipment.
\end{enumerate}
\textbf{Part IV.}
\begin{enumerate}[resume]
\item \emph{Named result}: in the model a 59-kilometre hop at \qty{11}{\giga\hertz} loses its \omterm{def:nw:long-haul-wireless:budget}{fade margin} at \qty{18.7}{\milli\metre\per\hour} of rain
  along its whole length (\qty{80.6}{\milli\metre\per\hour} in a ten-kilometre cell); a \qty{100}{\milli\metre\per\hour} storm crossing the route puts
  it on fibre for 21 minutes, \qty{2.93}{\milli\second} slower each way.
\item Lower adverse selection and lower trading costs while the fastest traders' microwave networks were disrupted.
\item Diversity against the weather: the lower band survives rain that stops the higher one, which carries more; together they keep more hours on
  the radio.
\item Detect the fade in milliseconds, move the latency-critical traffic to fibre, tell the strategies that their latency assumptions changed
  (widen, pause or stop the races they can no longer win), and move back when the hop recovers.
\item The storm climate along the route (rain-rate statistics) and the provider's real equipment figures and measured outages.
\item Its counterparties are slower for a while: quotes stay up longer and adverse selection falls; it may trade more cheaply.
\item Yes if the firm keeps trading during storms; a straighter fibre or hollow-core glass would narrow the gap.
\item It switches the fastest traders off for a few minutes or hours, and the market is measurably cheaper while it does.
\end{enumerate}
\end{solution}

\begin{solution}{iq:nw:long-haul-wireless:1}
Radio travels in air at nearly $c_0$ along an almost straight line of towers; light in glass travels at $c_0/1.46$ along the ground's rights of way.
The price is bandwidth (tens to hundreds of megabits against terabits), weather (\omterm{def:nw:long-haul-wireless:rain}{rain fade}) and towers to lease and maintain.
\iqlookfor{Medium and path; bandwidth and availability as costs; fibre as back-up.}
\end{solution}

\begin{solution}{iq:nw:long-haul-wireless:2}
The Earth's curvature and obstacles: the towers must lift the line of sight and the first \omterm{def:nw:long-haul-wireless:fresnel}{Fresnel zone} above the bulge, which grows with the square of
the hop; and the \omterm{def:nw:long-haul-wireless:budget}{link budget}: free-space loss and rain take the \omterm{def:nw:long-haul-wireless:budget}{fade margin} away as the hop lengthens.
\iqlookfor{Bulge, \omterm{def:nw:long-haul-wireless:fresnel}{Fresnel zone}, tower height, \omterm{def:nw:long-haul-wireless:budget}{link budget}; the 70-km order of magnitude.}
\end{solution}

\begin{solution}{iq:nw:long-haul-wireless:3}
The radio route may fade; traffic moves to fibre, three milliseconds slower. The systems must detect it, switch, adjust the strategies'
latency assumptions (fewer races, wider quotes, or pauses), and switch back; the firm's competitors on radio are in the same weather.
\iqlookfor{Detection and switch-over; strategy behaviour; everyone's radio fades together.}
\end{solution}

\begin{solution}{iq:nw:long-haul-wireless:4}
Treat storms along the route as exogenous shocks that remove the fastest traders' speed advantage; compare liquidity, adverse selection and
trading costs in those windows with similar windows without storms, controlling for weather's other effects; that is Shkilko and Sokolov's design.
\iqlookfor{Natural experiment; identification; controls for weather affecting demand itself.}
\end{solution}

\begin{solution}{iq:nw:long-haul-wireless:5}
The few fields that move the other market (best prices and trades of the leading instruments, or a derived signal), in a compact binary format,
and orders if the link carries them; everything else, including full depth and recovery, on fibre.
\iqlookfor{Bandwidth budget; message design; recovery path.}
\end{solution}

\begin{solution}{iq:nw:long-haul-wireless:6}
At 50~km: millimetre wave and lasers (gigabits, short hops, weather diversity). At 1\,000~km: a microwave chain in the 6--11-gigahertz bands with a
fibre back-up. At 6\,000~km across an ocean: a subsea cable for everything, and possibly shortwave for a few messages when the ionosphere allows.
\iqlookfor{Distance, bandwidth and weather as the axes; always a fibre back-up.}
\end{solution}
